\section{Discussion and Future Directions}
\label{sec:discussion}

This section explains the parameter choices, nearby extensions, and finite-dimensional effectiveness questions.

\subsection{Why degrees three and five}

The low-order inverse relations \eqref{eq:low-inverse} explain the parameter search.  The hyperplane scheme has $H(t)=\arcsin t$, so its inverse $\sin\zeta$ has alternating nonlinear coefficients.  To increase the admissible radius, the first objective is to suppress the cubic inverse obstruction.  An anti-aligned third-chaos average changes $b_3$ in the correct direction.  Once $b_3$ is controlled, the next obstruction is quintic, and an aligned fifth-chaos average gives the next degree of freedom.

The small two-dimensional perturbation $\eta \He_3(X)$ in the base threshold is also important.  A threshold depending only on the linear combination $Z+\sum_ds_dR_d$ would reduce to a hyperplane after rotation.  The $X$-dependent cubic term changes the coefficient integrals $\mathsf A_{a,b}$ and gives the higher coefficients the needed shape.

\subsection{A degree \texorpdfstring{$7/9$}{7/9} extension}

The coefficient formula immediately extends to additional odd degrees.  For instance, taking
\[
  D=\{3,5,7,9\},
  \qquad
  \sigma_d=(-1)^{(d-1)/2},
\]
keeps the same formula \eqref{eq:general-bm}.  A preliminary numerical search gives
\begin{align*}
  \eta&=0.12569605,&
  s_3&=0.34199778,&
  s_5&=0.05153822,\\
  s_7&=0.00698313,&
  s_9&=0.00090252.
\end{align*}
For the corresponding limiting model, the estimated value is
\[
  \gamma\approx0.88157525281,
  \qquad
  \KG\lesssim1.78180628572.
\]
This extension is not used in the main theorem.  It is included to indicate that the third/fifth scheme is not isolated: the same analytic and finite-transfer machinery applies to a broader family of Gaussianized Hermite-chaos averages.

\subsection{Effective dimensions}

Corollary~\ref{cor:finite-scheme} is qualitative: it proves the existence of sufficiently large finite length $N$, but it does not give a useful dimension bound.  For the Grothendieck constant as a mathematical constant, qualitative existence of a finite scheme is enough.  For an explicit algorithmic implementation, one should quantify the convergence $H_N\to H$ on the contour used in Theorem~\ref{thm:finite-transfer}.

The natural route is a collision expansion, by which we mean an expansion separating terms with distinct indices from terms in which some indices repeat.  For fixed degree $d$ and fixed $k$, the polynomial $\He_k(P_{d,N})$ has a leading distinct-index component
\[
  \sqrt{k!}\,N^{-k/2}
  \sum_{1\le i_1<\cdots<i_k\le N}
  \He_d(U_{d,i_1})\cdots \He_d(U_{d,i_k}),
\]
which is homogeneous of total Hermite degree $dk$ and has $L^2$-norm squared $(N)_k/N^k\to1$.  The remaining collision terms have vanishing $L^2$ mass for fixed $d,k$.  This explains at the coefficient level why a limiting Hermite mode of order $k$ contributes correlation power $t^{dk}$.  Making this expansion uniform to the required degree would produce explicit finite-dimensional estimates.
